\documentclass[a4paper,11pt]{article}
\usepackage[T1]{fontenc}
\usepackage[margin=25mm]{geometry}
\usepackage{array}
\usepackage{hyperref}
\usepackage[skipempty]{credits}

\title{The \texttt{credits} package}
\author{Bastian Rieck and Jannis Born}
\date{1 November 2026}

\begin{document}
\maketitle

The \texttt{credits} package produces contributor statements and heatmaps
using the Contributor Role Taxonomy (CRediT). It supports named roles and
numeric contribution values. The taxonomy is described at
\url{https://credit.niso.org/}.

\section{Installation}

Copy \texttt{credits.sty} into the directory containing your document.
Load it in the preamble:
\begin{verbatim}
\usepackage{credits}
\end{verbatim}
The package requires \texttt{ifthen}, \texttt{kvoptions}, \texttt{pgfkeys},
\texttt{tikz}, \texttt{xstring}, and \texttt{etoolbox}. These packages are
available in TeX Live and MiKTeX.

\section{Named roles}

Add each contributor with \verb|\credit{name}{roles}|. Separate roles with
commas. Their order is arbitrary, but their spelling and capitalization
must match the taxonomy in Section~\ref{sec:roles}. Write an ampersand as
\verb|\&|.

\begin{verbatim}
\credit{Marie Curie}{Investigation, Methodology}
\credit{Michael Faraday}{Investigation, Visualization}
\credit{Ada Lovelace}{Software, Writing -- original draft}

\insertcreditsstatement
\end{verbatim}

\credit{Marie Curie}{Investigation, Methodology}
\credit{Michael Faraday}{Investigation, Visualization}
\credit{Ada Lovelace}{Software, Writing -- original draft}
\noindent\insertcreditsstatement

\medskip
Use \verb|\insertcredits| to show the same contributions as a heatmap.
The command inserts a \texttt{tikzpicture}; it does not create a float.
Short contributor labels can make the heatmap easier to fit on a page:

\begin{verbatim}
\resetcredits
\credit{M.C.}{Investigation, Methodology}
\credit{M.F.}{Investigation, Visualization}
\credit{A.L.}{Software, Writing -- original draft}
\insertcredits
\end{verbatim}

\resetcredits
\credit{M.C.}{Investigation, Methodology}
\credit{M.F.}{Investigation, Visualization}
\credit{A.L.}{Software, Writing -- original draft}
\begin{center}
\insertcredits
\end{center}

\section{Numeric contributions}

Alternatively, give each contributor 14 numbers between 0 and 1.
The positions follow the taxonomy order in Section~\ref{sec:roles}.
Zero means no contribution; larger values produce darker cells.
Use the granular commands with numeric lists:

\begin{verbatim}
\resetcredits
\credit{M.C.}{1,0.5,0,0,1,1,0,0,0,0,0,0,0,0}
\credit{A.L.}{0,0,0,0,0,0,0,0,1,0,0,0.5,1,0}
\insertcreditsgranular
\insertcreditsgranularstatement
\end{verbatim}

\resetcredits
\credit{M.C.}{1,0.5,0,0,1,1,0,0,0,0,0,0,0,0}
\credit{A.L.}{0,0,0,0,0,0,0,0,1,0,0,0.5,1,0}
\begin{center}
\insertcreditsgranular
\end{center}
\noindent\insertcreditsgranularstatement

\medskip
The numeric text statement treats every nonzero value as a contribution;
it does not print the magnitude. Keep named roles and numeric lists in
separate blocks, using \verb|\resetcredits| between them.

\section{Options and commands}

Pass options when loading the package, for example:
\begin{verbatim}
\usepackage[role=blue,grid=gray,skipempty,horizontal]{credits}
\end{verbatim}

\noindent\begin{tabular}{@{}lp{0.72\linewidth}@{}}
\texttt{role} & Cell colour for contributions; default: black.\\
\texttt{grid} & Cell border colour; default: white.\\
\texttt{horizontal} & Put roles in columns and contributors in rows.\\
\texttt{skipempty} & Omit roles with no contributions. For named text grouped
  by author, omit contributors with no recognized roles.\\
\texttt{byauthor} & Group named text statements by author instead of role.\\
\texttt{separator} & Separate text entries with this value; default: a
  semicolon followed by a space.\\
\end{tabular}

\medskip
Boolean options are false unless specified. This manual loads
\texttt{skipempty}, so its heatmaps and statements show only used roles.
Define custom colours with \verb|\definecolor| before loading the package.
For example, \verb|separator={\newline}| puts each text entry on a new line.

\begin{description}
\item[\texttt{\textbackslash credit\{name\}\{roles or numbers\}}]
  Add a contributor. Unknown named roles and out-of-range numbers produce
  package warnings.
\item[\texttt{\textbackslash insertcredits}]
  Draw a heatmap from named roles.
\item[\texttt{\textbackslash insertcreditsstatement}]
  Print a text statement from named roles.
\item[\texttt{\textbackslash insertcreditsgranular}]
  Draw a heatmap from numeric contributions.
\item[\texttt{\textbackslash insertcreditsgranularstatement}]
  Print a text statement from numeric contributions, by role.
\item[\texttt{\textbackslash resetcredits}]
  Remove all stored contributors before starting another statement.
\end{description}

Names can contain escaped LaTeX characters such as \verb|\_| and
\verb|\$|. The package retains contributors until \verb|\resetcredits|
is called; printing a statement does not clear them.

\section{Taxonomy order}\label{sec:roles}

Named roles use the spelling below. Numeric lists use the same order.
\begin{enumerate}
\item Conceptualization
\item Data curation
\item Formal analysis
\item Funding acquisition
\item Investigation
\item Methodology
\item Project administration
\item Resources
\item Software
\item Supervision
\item Validation
\item Visualization
\item \verb|Writing -- original draft|
\item \verb|Writing -- review \& editing|
\end{enumerate}

\section{Examples and licence}

The distribution includes \texttt{example.tex} (named heatmap),
\texttt{example\_granular.tex} (numeric heatmap and text), and
\texttt{example\_text.tex} (numeric text). Compile them with
\texttt{pdflatex} from the directory containing \texttt{credits.sty}.

The package and this documentation are distributed under the BSD
3-Clause licence; see \texttt{LICENSE.md}. Source and issue tracker:
\url{https://github.com/Pseudomanifold/latex-credits}.

\end{document}
